\documentclass[11pt,a4paper]{article}

% ── Encoding & fonts ─────────────────────────────────────────────────────────
\usepackage[T1]{fontenc}
\usepackage[utf8]{inputenc}
\usepackage{lmodern}
\usepackage{microtype}

% ── Page geometry ─────────────────────────────────────────────────────────────
\usepackage[top=2.5cm, bottom=2.5cm, left=2.8cm, right=2.8cm]{geometry}

% ── Typography & spacing ──────────────────────────────────────────────────────
\usepackage{parskip}
\setlength{\parskip}{4pt}
\usepackage{setspace}
\setstretch{1.05}

% ── Colours ───────────────────────────────────────────────────────────────────
\usepackage[dvipsnames,table]{xcolor}
\definecolor{accent}{HTML}{1A56DB}
\definecolor{rowalt}{HTML}{F3F6FC}
\definecolor{codebackground}{HTML}{F5F5F5}
\definecolor{rulecolor}{HTML}{CCCCCC}

% ── Section headings ──────────────────────────────────────────────────────────
\usepackage{titlesec}
\titleformat{\section}{\large\bfseries\color{accent}}{}{0em}{}[\color{rulecolor}\titlerule]
\titleformat{\subsection}{\normalsize\bfseries}{}{0em}{}
\titleformat{\subsubsection}{\normalsize\itshape}{}{0em}{}
\titlespacing{\section}{0pt}{14pt}{4pt}
\titlespacing{\subsection}{0pt}{9pt}{3pt}
\titlespacing{\subsubsection}{0pt}{6pt}{2pt}

% ── Tables ────────────────────────────────────────────────────────────────────
\usepackage{booktabs}
\usepackage{tabularx}
\usepackage{array}
\renewcommand{\arraystretch}{1.22}
\newcolumntype{L}[1]{>{\raggedright\arraybackslash}p{#1}}
\newcolumntype{C}[1]{>{\centering\arraybackslash}p{#1}}

% ── Code listings ─────────────────────────────────────────────────────────────
\usepackage{listings}
\lstset{
  backgroundcolor=\color{codebackground},
  basicstyle=\ttfamily\small,
  breaklines=true,
  breakatwhitespace=true,
  frame=single,
  framerule=0.4pt,
  rulecolor=\color{rulecolor},
  xleftmargin=8pt,
  xrightmargin=8pt,
  aboveskip=6pt,
  belowskip=6pt,
  captionpos=b,
  showstringspaces=false,
}

% ── Hyperlinks ────────────────────────────────────────────────────────────────
\usepackage[hidelinks,
  pdftitle={Feature Implementation Reference -- MyPDF Flutter},
  pdfauthor={MyPDF Engineering Team},
  colorlinks=false
]{hyperref}

% ── Header / footer ───────────────────────────────────────────────────────────
\usepackage{fancyhdr}
\pagestyle{fancy}
\fancyhf{}
\renewcommand{\headrulewidth}{0.4pt}
\fancyhead[L]{\small\textcolor{gray}{MyPDF --- Feature Implementation Reference}}
\fancyhead[R]{\small\textcolor{gray}{28 May 2026}}
\fancyfoot[C]{\small\thepage}

% ── Misc ──────────────────────────────────────────────────────────────────────
\usepackage{enumitem}
\setlist[itemize]{itemsep=1pt, topsep=3pt, parsep=0pt}
\setlist[enumerate]{itemsep=1pt, topsep=3pt, parsep=0pt}

% ─────────────────────────────────────────────────────────────────────────────
\begin{document}

% ── Compact title block ───────────────────────────────────────────────────────
\thispagestyle{fancy}
\begin{center}
  {\LARGE\bfseries\color{accent} Feature Implementation Reference}\\[6pt]
  {\normalsize MyPDF \textbullet{} Flutter (Android + Web) \textbullet{}
   Firebase + Supabase \textbullet{} Riverpod \textbullet{} 28 May 2026}\\[4pt]
  {\small\itshape Prepared by the MyPDF Engineering Team \textbullet{}
   Private --- not for redistribution}
\end{center}
\vspace{2pt}
{\small This document describes the complete implementation logic for every feature in
the MyPDF Flutter application: authentication, navigation, PIN security, library
management, PDF pipeline, local caching, reader, notes, text-to-speech, OCR,
profile, feature flags, and the Supabase CORS proxy.}
\vspace{8pt}
\hrule
\vspace{4pt}

\tableofcontents
\vspace{6pt}
\hrule

% ─────────────────────────────────────────────────────────────────────────────
\section{Architecture Overview}

\subsection{Clean Architecture Layers}

Every feature under \texttt{lib/features/<name>/} is divided into three mandatory layers:

\begin{enumerate}
  \item \textbf{data/} --- SDK adapters: \texttt{FirebaseAuth}, \texttt{FirebaseFirestore},
    Supabase Storage, Hive, HTTP. This layer may import any package.
  \item \textbf{domain/} --- Plain Dart models and repository interfaces only.
    Zero imports of \texttt{flutter/}, \texttt{firebase\_*}, \texttt{supabase\_*},
    \texttt{hive*}, \texttt{path\_provider}, or \texttt{dart:io} are permitted here.
    Repository methods return \texttt{Either<Failure,\,T>} from \texttt{dartz}, using
    \texttt{lib/core/errors/failures.dart} (\texttt{AuthFailure}, \texttt{ServerFailure}).
  \item \textbf{presentation/} --- Riverpod controllers, screens, and widgets.
    Controllers are typed \texttt{AsyncNotifier<void>}; screens consume providers only
    via \texttt{ref.watch} and \texttt{ref.read}.
\end{enumerate}

The \texttt{reader/} and \texttt{profile/} features are \emph{presentation-only}: they
reuse \texttt{library/} data and domain layers rather than defining their own.

\subsection{Two-Backend Split}

\vspace{4pt}
\begin{tabularx}{\textwidth}{@{}L{2.8cm} L{11.0cm}@{}}
  \toprule
  \textbf{Backend} & \textbf{Responsibility} \\
  \midrule
  \rowcolor{rowalt}
  \textbf{Firebase} &
    Authentication (email/password) \textbullet{} Firestore document store
    (\texttt{users}, \texttt{bookshelves}, \texttt{books}, \texttt{notes}) \textbullet{}
    Crashlytics (mobile only) \textbullet{} Remote Config (feature flags) \\
  \textbf{Supabase} &
    Object Storage bucket \texttt{pdfs} (path \texttt{\{uid\}/\{millis\}.pdf}) \textbullet{}
    Edge Function \texttt{pdf-proxy} (Deno CORS bypass for web) \\
  \bottomrule
\end{tabularx}
\vspace{2pt}

Firebase Storage is \textbf{never used}. All PDF bytes live in the Supabase bucket.
All dates in Firestore are stored as \textbf{ISO 8601 strings} (not
\texttt{Timestamp}), keeping domain models free of Firebase SDK types.

\subsection{Platform Branching Discipline}

Four compile-time conditional imports isolate mobile-only packages from the Dart-to-JS
compiler:

\vspace{4pt}
\begin{tabularx}{\textwidth}{@{}L{5.0cm} L{3.8cm} L{4.8cm}@{}}
  \toprule
  \textbf{Interface file} & \textbf{Mobile impl} & \textbf{Web impl} \\
  \midrule
  \rowcolor{rowalt}
  \texttt{ocr\_data\_source.dart} & \texttt{tesseract\_ocr} FFI & Tesseract.js Web Worker \\
  \texttt{tts\_engine.dart} & \texttt{flutter\_tts} & W3C \texttt{SpeechSynthesis} \\
  \rowcolor{rowalt}
  \texttt{pdf\_text\_extractor.dart} & \texttt{flutter\_pdf\_text} & Syncfusion bytes \\
  \texttt{mobile\_pdf\_reader.dart} & \texttt{flutter\_pdfview} & stub (web uses \texttt{pdfx}) \\
  \bottomrule
\end{tabularx}
\vspace{2pt}

The conditional import selector is \texttt{dart.library.js\_interop}: if the library is
present the web implementation is chosen, otherwise the mobile one. Never import
\texttt{\_io.dart} or \texttt{\_web.dart} files directly; always depend on the interface.

\subsection{Web Layout}

\texttt{MyPdfApp} injects a responsive layout shell via
\texttt{MaterialApp.router(builder:)}:

\begin{itemize}
  \item \textbf{$<$ 600 px} --- pass-through (mobile form factor).
  \item \textbf{600--1023 px} --- \texttt{\_PhoneFrame}: centered 412$\times$896
    card with rounded corners and shadow, mimicking a phone in a browser window.
  \item \textbf{$\geq$ 1024 px} --- \texttt{\_DesktopRouteAwareShell}: listens to the
    router delegate; wraps authenticated routes in \texttt{DesktopShell}
    (sidebar + main content); shows \texttt{DesktopAuthShell} for login/register.
\end{itemize}

% ─────────────────────────────────────────────────────────────────────────────
\section{Initialisation Flow}

\textbf{Source:} \texttt{lib/main.dart}

The \texttt{main()} function runs the following sequence before
\texttt{runApp()} is called:

\begin{enumerate}
  \item \textbf{Hive open.} \texttt{Hive.openBox('app\_prefs')} makes the local cache
    available synchronously to every service that needs it.
  \item \textbf{Firebase init.} \texttt{Firebase.initializeApp()} using
    \texttt{DefaultFirebaseOptions.currentPlatform} (generated by FlutterFire CLI).
    Firestore offline persistence is enabled:
    \texttt{Settings(persistenceEnabled: true, cacheSizeBytes:\,CACHE\_SIZE\_UNLIMITED)}.
  \item \textbf{Supabase init.} \texttt{Supabase.initialize(url, anonKey)} using the
    public anon key (safe to ship; service keys never committed).
  \item \textbf{Remote Config fetch.} \texttt{FeatureFlags.initialize()} fetches with a
    10-second timeout and a 1-hour minimum cache TTL. Fetch failures are caught and
    logged; the app falls back to \texttt{\_defaults} without throwing.
  \item \textbf{\texttt{runApp}.} A \texttt{ProviderScope} with
    \texttt{overrides:\,[featureFlagsProvider.overrideWithValue(featureFlags)]} wraps
    \texttt{MyPdfApp}, injecting the pre-built \texttt{FeatureFlags} singleton so no
    async work happens inside providers.
\end{enumerate}

\subsection{Error Instrumentation}

Three Crashlytics capture sites are registered, all guarded by \texttt{if (!kIsWeb)}
and \texttt{setCrashlyticsCollectionEnabled(!kDebugMode)}:

\begin{itemize}
  \item \texttt{FlutterError.onError} $\rightarrow$ \texttt{recordFlutterFatalError}
    (widget build/render errors)
  \item \texttt{PlatformDispatcher.instance.onError} $\rightarrow$
    \texttt{recordError(fatal: true)} (platform-thread and isolate errors)
  \item \texttt{runZonedGuarded} catch $\rightarrow$ \texttt{recordError(fatal: true)}
    (all remaining async errors)
\end{itemize}

% ─────────────────────────────────────────────────────────────────────────────
\section{Authentication}

\subsection{Data Layer}

\textbf{Source:} \texttt{lib/features/auth/data/firebase\_auth\_data\_source.dart}

\texttt{FirebaseAuthDataSource} wraps \texttt{FirebaseAuth} and \texttt{FirebaseFirestore}:

\begin{itemize}
  \item \textbf{login(email, password):} Calls \texttt{signInWithEmailAndPassword},
    then \texttt{\_fetchUser(uid)} to return a \texttt{UserModel} with Firestore fields.
  \item \textbf{register(name, email, password):} Calls
    \texttt{createUserWithEmailAndPassword}, then writes
    \texttt{users/\{uid\}} with \texttt{\{name, email, role:\,'user',
    createdAt:\,serverTimestamp()\}}.
  \item \textbf{\_fetchUser(uid):} Fetches \texttt{users/\{uid\}}; falls back to
    \texttt{auth.currentUser.displayName/email} if the Firestore doc is missing (handles
    auth--Firestore creation race).
  \item \textbf{authStateChanges():} Streams \texttt{FirebaseAuth.authStateChanges()}
    and \texttt{asyncMap}-fetches the user profile for each non-null event.
\end{itemize}

\subsection{Error Mapping}

\textbf{Source:} \texttt{lib/features/auth/data/auth\_repository\_impl.dart}

All \texttt{FirebaseAuthException} codes are mapped to user-facing \texttt{AuthFailure}
messages. Login collapses \texttt{user-not-found}, \texttt{wrong-password}, and
\texttt{invalid-credential} into a single ``Incorrect email or password'' message to
prevent email-enumeration oracles. Other codes map to: \texttt{email-already-in-use},
\texttt{weak-password}, \texttt{too-many-requests}, \texttt{network-request-failed},
and \texttt{user-disabled}.

\subsection{Riverpod Providers}

\vspace{4pt}
\begin{tabularx}{\textwidth}{@{}L{4.8cm} L{2.8cm} L{6.0cm}@{}}
  \toprule
  \textbf{Provider} & \textbf{Type} & \textbf{Purpose} \\
  \midrule
  \rowcolor{rowalt}
  \texttt{authStateProvider} & \texttt{StreamProvider} (keepAlive) &
    \texttt{Stream<UserModel?>} fed by Firebase auth changes; watched by the router \\
  \texttt{authControllerProvider} & \texttt{AsyncNotifier} &
    State machine: idle $\to$ loading $\to$ error/success \\
  \rowcolor{rowalt}
  \texttt{userProfileProvider} & \texttt{StreamProvider} (keepAlive) &
    Direct Firestore \texttt{users/\{uid\}} watch for real-time name sync \\
  \bottomrule
\end{tabularx}

\subsection{Screens}

\texttt{LoginScreen} and \texttt{RegisterScreen} each expose a mobile form and a
desktop variant rendered inside \texttt{DesktopAuthShell} (a centred split-card layout).
Both listen to \texttt{authControllerProvider} and surface \texttt{Failure.message}
via a \texttt{SnackBar}. On success they call \texttt{context.go(AppRoutes.home)};
GoRouter's redirect logic takes over from there.

% ─────────────────────────────────────────────────────────────────────────────
\section{Navigation (GoRouter)}

\textbf{Sources:} \texttt{lib/core/constants/app\_router.dart},
\texttt{lib/core/constants/app\_routes.dart}

\subsection{Route Table}

\vspace{4pt}
\begin{tabularx}{\textwidth}{@{}L{4.8cm} L{8.8cm}@{}}
  \toprule
  \textbf{Route} & \textbf{Screen} \\
  \midrule
  \rowcolor{rowalt}
  \texttt{/login} & \texttt{LoginScreen} \\
  \texttt{/register} & \texttt{RegisterScreen} \\
  \rowcolor{rowalt}
  \texttt{/pin-setup} & \texttt{PinSetupScreen} (mandatory on first login) \\
  \texttt{/pin-entry} & \texttt{PinEntryScreen} (every cold start) \\
  \rowcolor{rowalt}
  \texttt{/home} & \texttt{HomeScreen} (shelves + recents) \\
  \texttt{/shelf/:id} & \texttt{ShelfContentScreen} \\
  \rowcolor{rowalt}
  \texttt{/book/new} & \texttt{NewBookScreen} (URL or upload) \\
  \texttt{/book/:id} & \texttt{BookInfoScreen} \\
  \rowcolor{rowalt}
  \texttt{/book/:id/reading} & \texttt{ReadingScreen} \\
  \texttt{/book/:id/note} & \texttt{NoteScreen} \\
  \rowcolor{rowalt}
  \texttt{/profile} & \texttt{ProfileScreen} \\
  \texttt{/profile/edit} & \texttt{EditProfileScreen} \\
  \bottomrule
\end{tabularx}

\subsection{Redirect Logic}

The pure function \texttt{computeRedirect(state, authState, hasPinSet, pinUnlocked)}
is the single decision point, evaluated on every navigation event:

\begin{enumerate}
  \item Auth state is loading $\Rightarrow$ return \texttt{null} (avoids a
    \texttt{/login} flash while Firebase resolves).
  \item User is not logged in and current location is not an auth route
    $\Rightarrow$ redirect to \texttt{/login}.
  \item User is logged in, PIN not yet set $\Rightarrow$ redirect to
    \texttt{/pin-setup}.
  \item User is logged in, PIN set but not unlocked for this session
    $\Rightarrow$ redirect to \texttt{/pin-entry}.
  \item All gates pass and user is on an auth or PIN route
    $\Rightarrow$ redirect to \texttt{/home} (bounce rule).
  \item Otherwise $\Rightarrow$ return \texttt{null} (allow navigation).
\end{enumerate}

\texttt{routerProvider} (keepAlive) watches \texttt{authStateProvider} and synchronously
reads \texttt{appPinServiceProvider.hasPinSet()} and \texttt{appPinSessionProvider} on
each redirect evaluation. Being a pure function, redirect logic is covered by 10 unit
tests in \texttt{test/core/router/router\_redirect\_test.dart} without spinning up a
GoRouter instance.

% ─────────────────────────────────────────────────────────────────────────────
\section{App-Level PIN and Biometric}

\subsection{PIN Storage}

\textbf{Source:} \texttt{lib/core/local/app\_pin\_service.dart}

The raw PIN is never stored. On \texttt{setPin(pin)}, the service generates a random
salt and stores only the \textbf{SHA-256-crypt} hash
(\texttt{\$5\$<salt>\$<hash>}) using \texttt{package:crypt}, written to Hive
\texttt{app\_prefs} under key \texttt{app\_pin\_hash}.

Verification calls \texttt{Crypt.match(pin, storedHash)}, which is constant-time to
prevent timing attacks. \texttt{clearPin()} deletes the key on logout.

\subsection{Session State}

\textbf{Source:} \texttt{lib/core/local/app\_pin\_session.dart}

\texttt{AppPinSession} is an in-memory Riverpod \texttt{StateNotifier<bool>} initialised
to \texttt{false}. \texttt{unlock()} sets state to \texttt{true};
\texttt{lock()} resets it. The state is \textbf{never persisted to disk} --- killing the
app re-locks it unconditionally.

\subsection{PIN Entry Screen}

\texttt{PinEntryScreen} renders a 6-digit numpad (3$\times$4 grid with 0 and
backspace). Key behaviours:

\begin{itemize}
  \item \textbf{Rate limiting.} Five consecutive wrong attempts trigger a 30-second
    cooldown. A countdown is shown inline; input is disabled during the window.
  \item \textbf{Biometric quick-unlock.} If \texttt{BiometricAuthService.isDeviceSupported()}
    returns \texttt{true} on mount, a ``Use Face ID / Fingerprint'' button is shown.
    Tapping it calls \texttt{authenticate(biometricOnly: true)} and unlocks the session
    on success.
  \item \textbf{Escape hatch.} A ``Forgot PIN? Sign out'' link calls
    \texttt{authControllerProvider.logout()}, clearing all state.
  \item \textbf{Keyboard support.} Numeric keys and backspace are handled via
    \texttt{LogicalKeyboardKey} mapping.
\end{itemize}

\subsection{PIN Setup Screen}

\texttt{PinSetupScreen} is mandatory on first login (no skip, no back button). It
operates in two steps: enter PIN, then confirm PIN. The screen auto-advances to the
confirm step when 6 digits are entered. A mismatch shows a shake animation and
``PINs don't match. Try again.'', then resets to step 1. On success, \texttt{setPin()}
is called, the session is marked unlocked, and the app navigates to \texttt{/home}.

\subsection{BiometricAuthService}

\textbf{Source:} \texttt{lib/features/auth/data/biometric\_auth\_service.dart}

A thin wrapper around \texttt{package:local\_auth} that provides:
\texttt{isDeviceSupported()} (hardware + enrollment check),
\texttt{authenticate()}, and a per-user opt-in flag
(\texttt{biometric\_enabled} in Hive \texttt{app\_prefs}). Web always returns
\texttt{false} from \texttt{isDeviceSupported()}. Biometric is an acceleration layer
only --- the PIN or email/password credential always remains valid.

% ─────────────────────────────────────────────────────────────────────────────
\section{Library --- Shelves and Books}

\subsection{Domain Models}

\vspace{4pt}
\begin{tabularx}{\textwidth}{@{}L{3.2cm} L{10.4cm}@{}}
  \toprule
  \textbf{Model} & \textbf{Fields} \\
  \midrule
  \rowcolor{rowalt}
  \texttt{BookshelfModel} &
    \texttt{id}, \texttt{name}, \texttt{ownerId}, \texttt{createdAt} (ISO 8601) \\
  \texttt{BookModel} &
    \texttt{id}, \texttt{title}, \texttt{link}, \texttt{totalPages},
    \texttt{currentPage}, \texttt{progress} (0--100), \texttt{status}
    (\texttt{reading}/\texttt{on\_hold}/\texttt{finished}), \texttt{shelfId},
    \texttt{ownerId}, \texttt{lastReadAt?}, \texttt{author?}, \texttt{year?},
    \texttt{needsOcr} (bool, immutable after creation) \\
  \rowcolor{rowalt}
  \texttt{NoteModel} &
    \texttt{id}, \texttt{bookId}, \texttt{title}, \texttt{content},
    \texttt{updatedAt} (ISO 8601) \\
  \bottomrule
\end{tabularx}
\vspace{2pt}

\texttt{BookModel.progress} is recomputed from
\texttt{currentPage / totalPages * 100} in \texttt{fromMap()}.
\texttt{link} is either an HTTPS URL (external or Supabase-hosted) or
\texttt{local://\{filename\}} (mobile-legacy only; rejected on web).

\subsection{Firestore Data Source}

\textbf{Source:} \texttt{lib/features/library/data/firestore\_data\_source.dart}

The data source exposes streams and mutation methods. Exceptions are not caught here;
they bubble up to the controller, which maps them to \texttt{AsyncError}.

Key stream methods: \texttt{watchShelves(ownerId)},
\texttt{watchBooks(ownerId)}, \texttt{watchBooksByShelf(shelfId, ownerId)},
\texttt{watchNotesByBookId(bookId)} (sorted \texttt{updatedAt} DESC),
\texttt{watchNotesByBookId}, \texttt{watchUserNotesCount(bookIds)}.

\texttt{updateReadingProgress()} skips \texttt{totalPages} and \texttt{progress} fields
if \texttt{totalPages == 0}, preventing corruption of pre-saved values when the page
count is not yet known.

\subsubsection{Cascading Delete Sequence}

Deleting a book involves five ordered steps:

\begin{enumerate}
  \item \texttt{FirestoreDataSource.deleteBook()} batch-deletes all associated notes
    (chunks of 500), then deletes the book document. Returns \texttt{book.link}.
  \item \texttt{RecentBooksService.remove(bookId)} --- removes from the LRU recents list.
  \item \texttt{OcrCacheService.purgeBook(bookId)} --- drops all
    \texttt{ocr\_v1\_\{bookId\}\_*} Hive entries (best-effort, exceptions logged).
  \item If \texttt{link} contains \texttt{/storage/v1/object/public/pdfs/},
    \texttt{Supabase.storage.from('pdfs').remove([path])} purges the object
    (best-effort, unawaited).
  \item On mobile, \texttt{pdf\_\{hash\}.pdf} and \texttt{thumb\_\{hash\}.jpg} are
    deleted from \texttt{appDocs} (best-effort).
\end{enumerate}

Deleting a shelf batch-updates all its books to \texttt{shelfId: ''}, then deletes the
shelf document. All batches are chunked at 500 operations (Firestore limit).

\subsection{whereIn Chunking}

Firestore \texttt{whereIn} accepts at most 30 values. \texttt{watchUserNotesCount(bookIds)}
splits \texttt{bookIds} into chunks of 30, spawns one Firestore listener per chunk, and
merges the per-chunk counts into a single \texttt{StreamController<int>.broadcast()}
stream. This powers the NOTES stat on the profile screen.

\subsection{Adding a Book}

\textbf{Source:} \texttt{lib/features/library/presentation/new\_book\_screen.dart}

The screen has two tabs:

\textbf{URL tab:}
\begin{enumerate}
  \item Validate URL scheme.
  \item Fetch bytes via \texttt{fetchPdfBytes(url)} (routes through CORS proxy on web
    for external hosts).
  \item \texttt{\_isBitmapOnlyPdf(bytes)}: samples first three pages via Syncfusion; if
    fewer than 5 characters of embedded text are found across all sampled pages, the PDF
    is flagged as bitmap-only and \texttt{needsOcr = true}. A dismissable warning is
    shown.
  \item \texttt{extractPdfMetadata(bytes)}: reads \texttt{author} and \texttt{year}
    from the PDF info dictionary (validated: author non-empty; year 1970--2099, rejecting
    the common default 2001-01-01).
  \item \texttt{primePdfCache(url, bytes)}: writes already-fetched bytes to the local
    cache path to avoid a re-download when the reader first opens.
  \item Duplicate-title guard: \texttt{\_norm(title)} (trim + lowercase) compared
    against all existing book titles.
  \item \texttt{LibraryController.createBook()} writes to Firestore.
\end{enumerate}

\textbf{Upload tab:} \texttt{FilePicker} is called with \texttt{withData: kIsWeb}
(bytes on web, path on mobile). After probing and metadata extraction, bytes are uploaded
to Supabase via \texttt{Supabase.storage.from('pdfs').uploadBinary()} with a random
suffix appended to the millisecond timestamp to guard against path collisions. The rest
of the pipeline is identical to the URL tab.

% ─────────────────────────────────────────────────────────────────────────────
\section{PDF Pipeline}

\subsection{PDF Fetcher}

\textbf{Source:} \texttt{lib/core/network/pdf\_fetcher.dart}

\texttt{fetchPdfBytes(url)} is the \textbf{single entry point} for all PDF network
requests. It must always be used; direct \texttt{http.get()} calls for PDF bytes are
prohibited.

\vspace{4pt}
\begin{tabularx}{\textwidth}{@{}L{4.0cm} L{9.6cm}@{}}
  \toprule
  \textbf{Platform / host} & \textbf{Request path} \\
  \midrule
  \rowcolor{rowalt}
  Mobile (any host) & Direct \texttt{http.get(url)}, 180 s timeout \\
  Web + Supabase host & Direct \texttt{http.get(url)} (same Supabase origin, no CORS
    issue) \\
  \rowcolor{rowalt}
  Web + external host & Proxied: \texttt{http.get(kCorsProxyBase + '?url=' +
    Uri.encodeComponent(url))} \\
  \bottomrule
\end{tabularx}
\vspace{2pt}

Non-200 responses throw an \texttt{Exception}. Proxy 502 errors include the first
200 characters of the response body for diagnosis. The constant \texttt{kCorsProxyBase}
holds the deployed Edge Function URL.

\subsection{pdfPathProvider}

\texttt{pdfPathProvider(url)} is a keepAlive \texttt{FutureProvider.family}:

\begin{itemize}
  \item \textbf{Web:} Returns \texttt{url} as-is. No filesystem access; \texttt{pdfx}
    loads from URL directly.
  \item \textbf{Mobile:} Checks for a cached file at
    \texttt{appDocs/pdf\_\{url.hashCode\}.pdf}. If absent or empty, downloads via
    \texttt{fetchPdfBytes()}, then validates:
    \begin{itemize}
      \item \texttt{\_looksLikePdf(bytes)}: scans the first 1\,100 bytes for
        \texttt{\%PDF-} signature.
      \item \texttt{\_hasEofMarker(bytes)}: scans the final 1\,024 bytes for
        \texttt{\%\%EOF} trailer.
    \end{itemize}
    Both checks must pass; otherwise an exception is thrown and the cached file is not
    written. This guards against truncated downloads and HTML error pages.
\end{itemize}

\subsection{Render Backends}

\vspace{4pt}
\begin{tabularx}{\textwidth}{@{}L{2.4cm} L{3.4cm} L{7.8cm}@{}}
  \toprule
  \textbf{Platform} & \textbf{Library} & \textbf{Notes} \\
  \midrule
  \rowcolor{rowalt}
  Mobile & \texttt{flutter\_pdfview} & Native Android/iOS PDFView; page callbacks via
    \texttt{onPageChanged(page, total)} \\
  Web & \texttt{pdfx} & Renders pages to \texttt{Uint8List} images; used for both the
    main reader view and all thumbnail generation \\
  \bottomrule
\end{tabularx}

\subsection{Metadata and Thumbnails}

\texttt{extractPdfMetadata(bytes)} (\texttt{lib/features/library/data/pdf\_metadata.dart})
uses \texttt{syncfusion\_flutter\_pdf} to read the document info dictionary. It never
throws; an unreadable PDF returns \texttt{PdfMetadata(author: null, year: null)}.

\texttt{pdfPageImageProvider(url, pageIndex)} renders a page to a 1600\,px-maximum JPEG
and is disk-cached on mobile. \texttt{pdfThumbnailProvider(url)} is an alias that
always requests page 0.

% ─────────────────────────────────────────────────────────────────────────────
\section{Local Cache (Hive)}

\textbf{Box:} \texttt{app\_prefs}, opened synchronously in \texttt{main()} before
\texttt{runApp()} is called.

\vspace{4pt}
\begin{tabularx}{\textwidth}{@{}L{4.0cm} L{3.4cm} L{6.3cm}@{}}
  \toprule
  \textbf{Service} & \textbf{Hive key(s)} & \textbf{Semantics} \\
  \midrule
  \rowcolor{rowalt}
  \texttt{RecentBooksService} & \texttt{recent\_book\_ids} &
    \texttt{List<String>}, max 10, LRU (insert-at-front + dedupe).
    Reactive via \texttt{\_box.watch(key:)} \\
  \texttt{OcrCacheService} & \texttt{ocr\_v1\_\{bookId\}\_\{pageIndex\}} &
    Per-page OCR text. Schema-versioned for engine swaps.
    Write failures are logged, never fatal \\
  \rowcolor{rowalt}
  \texttt{AppPinService} & \texttt{app\_pin\_hash} &
    SHA-256-crypt hash string \\
  \texttt{BiometricAuthService} & \texttt{biometric\_enabled} &
    Per-user opt-in boolean \\
  \bottomrule
\end{tabularx}
\vspace{2pt}

\textbf{Rule:} Only these four services write to \texttt{app\_prefs}. A new Hive key
requires a new service in \texttt{lib/core/local/}. No screen or provider may write to
the box directly.

% ─────────────────────────────────────────────────────────────────────────────
\section{Reader Feature}

\textbf{Source:} \texttt{lib/features/reader/presentation/reading\_screen.dart}

\subsection{Progress Tracking}

When the user turns a page, \texttt{onPageChanged(page, total)} fires (via
\texttt{flutter\_pdfview} on mobile; via viewport scroll detection on web). The handler
calls \texttt{LibraryController.updateProgress()}, which writes
\texttt{currentPage}, \texttt{totalPages}, \texttt{progress}, and
\texttt{lastReadAt} to Firestore. If \texttt{page == totalPages - 1} and
\texttt{book.status != 'finished'}, status is updated to \texttt{'finished'} automatically.

\subsection{Recent Books}

On cold-start open of a book, \texttt{RecentBooksService.markOpened(bookId)} is called
via \texttt{ref.read(recentBooksServiceProvider)}. This moves the book to position 0 in
the LRU list, capped at 10 items.

\subsection{PDF Byte Caching on Web}

On web, the reader caches \texttt{\_webPdfBytes} (the raw \texttt{Uint8List}) per
\texttt{book.link} key. This avoids a re-download on every TTS \texttt{\_speakCurrentPage}
call, since Syncfusion's text extractor requires bytes rather than a URL.

% ─────────────────────────────────────────────────────────────────────────────
\section{Notes Feature}

\subsection{Data Flow}

\textbf{Sources:} \texttt{lib/features/library/domain/note\_model.dart},
\texttt{lib/features/library/data/firestore\_data\_source.dart}

Notes live in a flat Firestore collection \texttt{notes/\{noteId\}} with fields
\texttt{bookId}, \texttt{title}, \texttt{content}, \texttt{createdAt},
\texttt{updatedAt}. Ownership resolves transitively: the Firestore rule for a note
fetches the parent book's \texttt{ownerId} via the \texttt{bookOwner()} helper.

CRUD methods: \texttt{createNote()}, \texttt{updateNote()}, \texttt{deleteNote()},
\texttt{deleteNotes(List<String> noteIds)} (chunked at 500).

\subsection{Providers and Screens}

\texttt{notesByBookProvider(bookId)} streams notes sorted by \texttt{updatedAt} DESC.
\texttt{userNotesCountProvider} aggregates note counts across all user books using the
30-item \texttt{whereIn} chunking pattern described in §6.

Auto-naming: when the user saves a note with a blank title, the screen sets the title
to ``Note (N)'' where N is the current count of notes for that book.

The desktop reader integrates notes in a right-side panel (MY NOTE tab) with an inline
editor; the mobile reader routes to \texttt{/book/:id/note}.

% ─────────────────────────────────────────────────────────────────────────────
\section{Text-to-Speech (TTS)}

\subsection{Engine Abstraction}

\textbf{Source:} \texttt{lib/features/reader/data/tts\_engine.dart}

\texttt{TtsEngine} is a platform-agnostic interface with a factory constructor that
resolves at compile time:

\vspace{4pt}
\begin{tabularx}{\textwidth}{@{}L{2.2cm} L{3.6cm} L{7.9cm}@{}}
  \toprule
  \textbf{Platform} & \textbf{Impl} & \textbf{Key details} \\
  \midrule
  \rowcolor{rowalt}
  Mobile & \texttt{tts\_engine\_io.dart} &
    \texttt{flutter\_tts} 4.2.0; \texttt{awaitSpeakCompletion(true)} prevents race
    with auto-advance; \texttt{setProgressHandler} fires per-word boundary events;
    ``interrupted'' and ``canceled'' errors are filtered \\
  Web & \texttt{tts\_engine\_web.dart} &
    W3C \texttt{window.speechSynthesis}; Completer-per-utterance prevents late
    \texttt{onerror} completing the wrong future; 120\,ms cancel delay before next
    \texttt{speak()} to let the ``interrupted'' event propagate \\
  \bottomrule
\end{tabularx}

\subsection{Voice Selection on Web}

\texttt{\_trySetWebVoice()} scores available voices (prefers Aria, Jenny, Jane, Libby,
Michelle) and selects the best match. Chrome populates the voice roster asynchronously
after the first user gesture. If the roster is empty on mount, \texttt{\_trySetWebVoice}
is retried just-in-time before each \texttt{speak()} call. Firefox does not always
provide a \texttt{charLength} in word-boundary events; a forward-whitespace-scan
workaround fills in the missing value.

\subsection{Text Extraction}

\vspace{4pt}
\begin{tabularx}{\textwidth}{@{}L{2.2cm} L{4.2cm} L{7.3cm}@{}}
  \toprule
  \textbf{Platform} & \textbf{Source} & \textbf{Indexing} \\
  \midrule
  \rowcolor{rowalt}
  Mobile & \texttt{flutter\_pdf\_text} & 1-indexed pages; \texttt{PdfDoc} handle
    cached across calls \\
  Web & Syncfusion \texttt{PdfTextExtractor} & 0-indexed
    \texttt{startPageIndex/endPageIndex}; bytes cached per \texttt{book.link} \\
  \bottomrule
\end{tabularx}
\vspace{2pt}

All extracted text passes through \texttt{tts\_text\_cleaner.dart}, which normalises
whitespace, removes soft hyphens, and strips diacritics.

\subsection{Karaoke Word Highlighting}

\textbf{Source:} \texttt{lib/features/reader/presentation/controllers/karaoke\_controller.dart}

\texttt{KaraokeController} maintains \texttt{fullText}, \texttt{currentStart/End}
(character offsets), \texttt{isSpeaking}, and \texttt{fallbackSentenceMode}. The reading
screen calls \texttt{onTtsStart(fullText, baseOffset)} before each \texttt{speak()},
and \texttt{onProgress(text, start, end, word)} on each engine boundary event.

\textbf{Fallback mode:} If no progress event fires within 2 seconds, the controller
switches to sentence-level highlights generated by splitting \texttt{fullText} on
\texttt{(?<=[.\,!\,?])\textbackslash{}s+}. Each sentence becomes a separate utterance;
completion advances the sentence queue. Fallback mode is never re-armed once a real
boundary event has arrived (\texttt{\_karaokeProgressEverFired} flag).

\textbf{Click-to-seek:} Tapping a word in \texttt{KaraokeTextPane} calls
\texttt{\_seekTtsTo(wordStart)}, which slices the remaining text from \texttt{wordStart},
cancels the current utterance, and re-speaks from the new position with
\texttt{baseOffset} updated accordingly.

% ─────────────────────────────────────────────────────────────────────────────
\section{OCR Feature}

\subsection{Conditional Import Pattern}

\textbf{Source:} \texttt{lib/features/library/data/ocr\_data\_source.dart}

The interface declares \texttt{recognize(Uint8List bytes, \{String langs\})} and
\texttt{dispose()}. A Dart conditional import selects the implementation at compile time:

\begin{lstlisting}[language=Dart]
// ocr_data_source.dart
export 'ocr_data_source_io.dart'
    if (dart.library.js_interop) 'ocr_data_source_web.dart';
\end{lstlisting}

Never import either \texttt{\_io} or \texttt{\_web} file directly.

\subsection{Mobile Implementation}

\textbf{Source:} \texttt{lib/features/library/data/ocr\_data\_source\_io.dart}

\begin{enumerate}
  \item Write \texttt{imageBytes} (PNG) to a temp file (lossless for OCR accuracy).
  \item For payloads $>$ 2\,MB, use \texttt{compute()} to write off-isolate and avoid
    UI-thread hitch.
  \item Call \texttt{TesseractOcr.extractText(tempFile.path,
    OCRConfig(language: 'eng+tha', engine: OCREngine.tesseract))}.
  \item Delete the temp file in a \texttt{finally} block.
\end{enumerate}

Trained data (\texttt{tessdata\_best/eng.traineddata},
\texttt{tha.traineddata}) must be placed in \texttt{assets/tessdata/} before build.
The plugin copies them to \texttt{appDocs/tessdata/} on first call.

\subsection{Web Implementation}

\textbf{Source:} \texttt{lib/features/library/data/ocr\_data\_source\_web.dart}

\begin{enumerate}
  \item Lazy-load \texttt{web/ocr/tesseract.min.js} via a \texttt{<script>} tag on the
    first \texttt{recognize()} call; result cached globally.
  \item Spawn a Tesseract.js v5 Web Worker once per datasource lifetime
    (\textasciitilde{}5 s cost). Worker paths:
    \texttt{workerPath: 'ocr/worker.min.js'},
    \texttt{corePath: '/ocr/tesseract-core-simd-lstm.wasm.js'},
    \texttt{langPath: '/ocr/lang/'}.
    OEM 1 (LSTM-only), PSM 6 (single text block).
  \item Convert \texttt{Uint8List} $\rightarrow$ \texttt{Blob} $\rightarrow$ object URL.
  \item Call \texttt{worker.recognize(blobUrl)}; revoke the Blob URL in a
    \texttt{finally} block.
  \item On datasource disposal, call \texttt{worker.terminate()} (best-effort).
\end{enumerate}

\subsection{OCR Provider Pipeline}

\textbf{Source:} \texttt{lib/features/library/presentation/library\_providers.dart}
(\texttt{ocrPageTextProvider})

\begin{enumerate}
  \item \textbf{Cache check.} \texttt{OcrCacheService.get(bookId, pageIndex)} returns
    cached text synchronously if present (even an empty string means OCR was already
    attempted).
  \item \textbf{Image render.} \texttt{ocrPageImageProvider(url, pageIndex)} returns a
    2\,400\,px-maximum lossless PNG (\textasciitilde{}280\,DPI, within Tesseract's
    recommended 250--400\,DPI range).
  \item \textbf{Recognition.} \texttt{ocrDataSource.recognize(bytes, langs: 'eng+tha')}.
  \item \textbf{Text cleaning.} \texttt{cleanForTts()} normalises whitespace and removes
    artefacts.
  \item \textbf{Cache write.} \texttt{OcrCacheService.put(bookId, pageIndex, text)}
    (best-effort; failures are logged).
  \item \textbf{Invalidate rendered PNG.} \texttt{ref.invalidate(ocrPageImageProvider(...))}
    releases the lossless PNG from Riverpod's cache, critical for memory during long
    background sweeps.
\end{enumerate}

\subsection{Background Pre-OCR Sweep}

After the first successful foreground OCR recognition, \texttt{\_maybeStartBackgroundOcr}
spawns a loop that iterates all pages except the triggering page and any already-cached
pages, calling \texttt{ref.read(ocrPageTextProvider(...).future)} sequentially. A
version token \texttt{\_bgOcrVersion} is incremented on dispose or book-switch; the
loop checks the token at each iteration start to cancel cleanly. Progress is reported
to \texttt{bookOcrProgressProvider} \texttt{({done,\,total})?}, which drives a chip in
the reader app bar.

\subsection{Kill Switch}

\texttt{featureFlagsProvider.ocrFallbackEnabled} (backed by Remote Config key
\texttt{ocr\_fallback\_enabled}, default \texttt{true}) must be \texttt{true} for any
OCR path to execute. When \texttt{false}, the reading screen shows a snackbar
``No readable text (scanned PDF?)'' instead of invoking Tesseract.

% ─────────────────────────────────────────────────────────────────────────────
\section{Profile Feature}

\textbf{Sources:} \texttt{lib/features/profile/presentation/}

\texttt{ProfileScreen} watches \texttt{userProfileProvider} (real-time Firestore stream)
for the user's \texttt{name} and displays \texttt{email} from
\texttt{authStateProvider}. Three stat cards show READ (total book count via
\texttt{allBooksProvider}), NOTES (\texttt{userNotesCountProvider}), and SHELVES
(\texttt{shelvesProvider}).

\texttt{EditProfileScreen} presents a single editable field: \texttt{name} (email is
read-only from Firebase Auth). On save, \texttt{firestoreDataSource.updateUserProfile(uid, name:)} writes the new name to \texttt{users/\{uid\}.name}. There is no avatar field and
no profile-level biometric toggle.

% ─────────────────────────────────────────────────────────────────────────────
\section{Feature Flag System}

\textbf{Source:} \texttt{lib/core/config/feature\_flags.dart}

\texttt{FeatureFlags} wraps \texttt{FirebaseRemoteConfig} with typed getters so app code
never calls \texttt{FirebaseRemoteConfig.instance} directly.

\begin{lstlisting}[language=Dart]
bool get ocrFallbackEnabled =>
    _remoteConfig.getBool('ocr_fallback_enabled');
\end{lstlisting}

Initialization is \textbf{failure-safe}: a \texttt{try/catch} around \texttt{fetch} and
\texttt{activate} catches network errors or quota-exceeded exceptions, logs them via
\texttt{AppLogger.warn}, and falls through to the \texttt{\_defaults} map without
rethrowing. This ensures a cold-start offline device still launches with sensible
defaults.

\vspace{4pt}
\begin{tabularx}{\textwidth}{@{}L{4.4cm} C{1.8cm} L{7.4cm}@{}}
  \toprule
  \textbf{Remote Config key} & \textbf{Default} & \textbf{Effect} \\
  \midrule
  \rowcolor{rowalt}
  \texttt{ocr\_fallback\_enabled} & \texttt{true} &
    Master switch for the entire Tesseract OCR pipeline. Set to \texttt{false} to
    disable OCR across all clients within the 1-hour TTL, requiring zero redeployment. \\
  \bottomrule
\end{tabularx}

% ─────────────────────────────────────────────────────────────────────────────
\section{Supabase Edge Function --- CORS Proxy}

\textbf{Source:} \texttt{supabase/functions/pdf-proxy/index.ts} (Deno runtime)

\subsection{Purpose}

Flutter Web cannot fetch arbitrary cross-origin PDFs due to browser CORS policy. The
\texttt{pdf-proxy} Edge Function runs inside the Supabase project and adds permissive
CORS headers to the upstream response. It is deployed with
\texttt{--no-verify-jwt} to allow unauthenticated browser requests; SSRF defences
compensate for the open surface.

\subsection{Request Flow}

\begin{lstlisting}
Browser  GET /functions/v1/pdf-proxy?url=<encoded-pdf-url>
            |
            v
        ssrfCheck(hostname)  <-- block private IPs, metadata hosts
            |
            v
        fetch(url, {redirect: 'manual', signal: AbortSignal.timeout(120s)})
            |
            v
        validate Content-Type (must be application/pdf or octet-stream)
        validate Content-Length <= 50 MB
        streaming cap via TransformStream (abort at 50 MB)
            |
            v
        Response {
            Content-Type: application/pdf,
            Access-Control-Allow-Origin: *,
            Cache-Control: public, max-age=60, must-revalidate
        }
\end{lstlisting}

\subsection{SSRF Defences}

\begin{itemize}
  \item \textbf{Private IP rejection.} DNS A/AAAA records are resolved and checked
    against all RFC 1918 ranges (\texttt{10/8}, \texttt{172.16/12},
    \texttt{192.168/16}), loopback (\texttt{127/8}), link-local (\texttt{169.254/16}),
    and IPv6 ULA. \texttt{localhost}, \texttt{*.local}, \texttt{*.internal}, and
    \texttt{metadata.*} are blocked by name before any DNS lookup.
  \item \textbf{Dual-resolver fallback.} System DNS is tried first; for TLDs that may
    fail (e.g., \texttt{.cr}), Cloudflare DoH (\texttt{1.1.1.1}) is used as a fallback.
  \item \textbf{Post-fetch re-validation.} The resolved IP is re-checked after the fetch
    response arrives, blocking DNS-rebinding attacks where the initial resolution returns
    a public IP but a subsequent lookup returns an internal one.
  \item \textbf{Redirect rejection.} \texttt{fetch(\{redirect: 'manual'\})} prevents
    open-redirect chains from routing through attacker-controlled servers. Any redirect
    returns HTTP 502 to the client.
\end{itemize}

\end{document}
